\documentclass[10pt,a4paper]{amsart}
\usepackage[margin=19mm]{geometry}
\usepackage{amsmath,amssymb,mathtools}
\usepackage[scr=rsfs,cal=euler]{mathalfa}
\usepackage{xfrac}
\usepackage[unicode,hidelinks,bookmarks=false]{hyperref}
\newcommand{\ie}{i.e.\ }
\newcommand{\cf}{cf.\ }
\newcommand{\resp}{respectively\ }
\newcommand{\Subsection}{\subsection}
\providecommand{\nobreakdash}{\nobreak\hbox{-}}
\setlength{\emergencystretch}{2em}
\multlinegap0pt
\usepackage{amssymb}
\usepackage{bm}
\usepackage{mathtools}
\usepackage{xcolor}
\usepackage{enumerate}
\newtheorem{lem}{Lemma}
\newtheorem{prop}{Proposition}
\newcommand{\CI}{\mathrm{(CI)}} 
\newcommand{\Irr}{\mathrm {Irr}}
\renewcommand{\leq}{\leqs}
\newcommand{\leqs}{\leqslant}
\title{Some generalizations of Camina Pairs and orders of elements in cosets}
\author{Thu T.H. Quan, Hung P. Tong-Viet}
\date{Source: arXiv:2508.13056v2 (CC BY 4.0)}
\begin{document}
\maketitle

\noindent\emph{Source and licence.} Some generalizations of Camina Pairs and orders of elements in cosets. Version arXiv:2508.13056v2 (CC BY 4.0), distributed by arXiv under the Creative Commons Attribution 4.0 International licence, http://creativecommons.org/licenses/by/4.0/, as recorded in the arXiv licence metadata for this exact version and shown on the abstract page https://arxiv.org/abs/2508.13056v2. Full licence notice: \texttt{licenses/arxiv-2508.13056-CC-BY-4.0.txt}.\par\smallskip

\noindent\textit{Editorial scope.} Counted unit: the article's prop prop7 (source0.tex lines 397--399). The article's complete original proof of it is delivered with it (source0.tex lines 400--427, one \texttt{proof} environment of 28 lines). No mathematics is written, repaired or re-derived: the statement and the proof are the article's own lines, in the article's own notation, and no retained line is substituted or re-flowed.  Retained uncounted as the source-local context the proof needs: lines 243--245; lines 352--354.  Nothing was omitted from any retained span, so the package carries no omission record.  No retained line contains a citation, so the wrapper carries no bibliography and no bibliography entry.  No retained line contains a figure environment, an image-inclusion command or drawing code, so no image is delivered and the package declares no asset dependency.  Editorial formatting only, in addition: an AMS \texttt{amsart} preamble that restates the article's own declarations and macros used by the retained lines, namely the environments \texttt{lem}, \texttt{prop} on separate counters, together with the standard packages those lines need.  The retained environments print in this document as Lemma 1, Proposition 1 in order of appearance, whereas the article prints the same items with the numbering of its own full document.  The complete native source of the article as distributed in the arXiv e-print for this exact version is bundled unchanged as \texttt{sources/183455/source0.tex}.
\begin{lem} \label{Hnormal}
Let $G$ be a finite group and let $H$ be a subgroup of $G$.	Suppose that the pair $(G,H)$ satisfies condition $\CI$ and $H$ is not normal in $G$. Let $N$ be the normal closure of $H$ in $G$. Then  $H$ is a normal subgroup of $N$ and $N/H$ is abelian. 
\end{lem}

\begin{prop}\label{prop6}
Let $G$ be a finite group and let $H$ be a subgroup of $G$.	Suppose that the pair $(G,H)$ satisfies condition $\CI$ and that $H$ is not normal in $G$. Let $N$ be the normal closure of $H$ in $G$. Assume $\chi\in \Irr(G\mid N)$ such that $[\chi_H,1_H] \neq 0$. Then $\chi(xh) = \chi(x)$ for all $x\in N\setminus H$, $h\in H$.
\end{prop}

\begin{prop}\label{prop7}
Let $G$ be a finite group and let $H$ be a subgroup of $G$.	Suppose that the pair $(G,H)$ satisfies condition $\CI$ and that $H$ is not normal in $G$. Let $N$ be the normal closure of $H$ in $G$. Assume $\chi \in \Irr(G\mid N)$ such that $[\chi_H,1_H]=0$. Then $\chi(xh) = \chi(x)=0$ for all $x\in N\setminus H$ and $h\in H$.
\end{prop}

\begin{proof}
	First, by applying Clifford's Theorem with $N\unlhd G$, we obtain 
	\begin{align}\label{9}
		\chi_N = f(\theta_1+\theta_2+\ldots + \theta_t),
	\end{align}
	where $\theta_i$ for $i=1,2,\ldots,t$ are distinct irreducible characters of $N$, which are conjugate in $G$. Since $[\chi_H,1_H]=0$, we can find a nontrivial irreducible constituent $\mu_1$ of $\chi_H$.
	We may assume that $\theta_1,\ldots,\theta_{l_1}$, for some $l_1\leq t$, are precisely the irreducible constituents of $\chi_N$ lying over $\mu_1$. Analogous to the argument in Proposition \ref{prop6}, we claim $$\Irr(N\mid \mu_1) = \left\{\theta_1,\ldots,\theta_{l_1}\right\}\text{ and }
		\mu_1^N = a_1(\theta_1+\ldots+\theta_{l_1})$$
	for some positive integer $a_1$. To verify this, take a character $\varphi \in \left\{\theta_1,\ldots,\theta_{l_1}\right\}$. Then $\varphi$ is an irreducible character of $N$ lying over $\mu_1$, so $\varphi\in \Irr(N\mid \mu_1)$. Conversely, let $\theta\in \Irr(N\mid \mu_1)$. Since $(G,H)$ satisfies condition $\CI$ and $\mu_1\neq 1_H$, $\mu_1^G = (\mu_1^N)^G = e\chi$ for some positive integer $e$. Thus $\chi$ is the unique irreducible constituent of $\theta^G$, which implies that $\theta$ is an irreducible constituent of $\chi_N$. As $\theta$ lies over $\mu_1$, it must be among $ \left\{\theta_1,\ldots,\theta_{l_1}\right\}$. This shows that $\left\{\theta_1,\ldots,\theta_{l_1}\right\} = \Irr(N\mid \mu_1)$, as claimed.
	
	Next, since $H\unlhd N$ by Lemma \ref{Hnormal}, we assume that for some elements $n_1,\ldots,n_s\in N$,  the characters $\mu_1,\mu_1^{n_1},\ldots,\mu_1^{n_s}\in \Irr(H)$ are all the distinct $N$-conjugates of $\mu_1$.
	Because $\theta_i$ lies over $\mu_1$ for all $i=1,\ldots,l_1$, applying Clifford's Theorem for $H\unlhd N$, we obtain
	$$(\theta_i)_H = b_i(\mu_1+\mu_1^{n_1} + \ldots + \mu_1^{n_s}), $$
	for some positive integer $b_i$ and for all $i = 1,\ldots, l_1$. Since $\theta_1,\ldots,\theta_{l_1}$ have the same degree and since $\mu_1,\mu_1^{n_1},\ldots,\mu_1^{n_s}$ have the same degree,  it follows that
	$$	  b_i(s+1)\mu_1(1)= \theta_i(1)= \theta_1(1) = b_1 (s+1)\mu_1(1),$$  which implies that $b_1=b_2=\dots=b_{l_1}$. Denote this common value by $a_1$, we conclude that $a_1=[(\theta_i)_H,\mu_1] = [\theta_i,\mu_1^N]$ for all $i  = 1,\ldots,l_1$. Therefore, since $\Irr(N\mid\mu_1) = \left\{\theta_1,\ldots,\theta_{l_1}\right\}$,  it follows that 
	\[ \mu_1^N = a_1(\theta_1+\ldots+\theta_{l_1}) \text{ or equivalently } \theta_1 + \ldots + \theta_{l_1} = \dfrac{1}{a_1} \mu_1^N. \]
	
	We now proceed, as in the proof of  Proposition \ref{prop6}, to partition   the set $\left\{\theta_1,\theta_2,\ldots,\theta_t\right\}$ of irreducible constituents  of $\chi_N$. To that end, suppose that  $\mu_2,\ldots,\mu_n$ together with $\mu_1$ are all the irreducible constituents of $\chi_H$ that are not conjugate in $N$. We claim that  the sets $\Irr(N\mid \mu_1),\Irr(N\mid\mu_2),\ldots,\Irr(N\mid \mu_n)$ form a partition of $\left\{\theta_1,\ldots,\theta_t\right\}$. To see that these sets are pairwise disjoint, suppose otherwise. Then for some $j\neq k$, $\Irr(N\mid\mu_j)$ and $\Irr(N\mid \mu_k)$ would share a common element $\lambda\in\Irr(N)$, implying that both $\mu_j$ and $\mu_k$ appear in $\lambda_H$. However, since $H\unlhd N$, Clifford's Theorem would imply that $\mu_j$ and $\mu_k$ are conjugate in $N$, contradicting the assumption. 
	
	Now, since $\mu_j$ for $j=2,\ldots,n$ is a constituent of $\chi_H$, it must  lie under some $\theta_i\in \{\theta_1,\dots,\theta_t\}$. We may assume  that each $\mu_j$ lies under $\theta_{l_{j-1}+1},\ldots,\theta_{l_{j}}$ for $j=2,\ldots,n$ with $l_0=0$ and ${l_n} = t$. Repeating the previous argument, we conclude that for all $j=2,\ldots,n$, $$\left\{\theta_{l_{j-1}+1},\ldots,\theta_{l_j}\right\} = \Irr(N\mid \mu_j)\text{ and }
	\theta_{l_{j-1}+1} + \ldots + \theta_{l_j} = \dfrac{1}{a_j}\mu_j^N,$$
	for some positive integer $a_j$. Therefore, we obtain the desired partition $$\Irr(N\mid \mu_1)\cup \Irr(N\mid \mu_2)\cup \ldots\cup\Irr(N\mid\mu_n)=\left\{\theta_1,\ldots,\theta_t\right\}.$$ From Equation \eqref{9}, 
	 	\begin{align}\label{1}
		\chi_N &= f (\theta_1+ \ldots +\theta_{l_1} + \theta_{l_1+1}+\ldots + \theta_{l_2} + \ldots + \theta_{l_{n-1}+1} + \ldots +\theta_{l_n}) \nonumber\\
		& = f\left(\dfrac{1}{a_1}\mu_1^N + \dfrac{1}{a_2}\mu_2^N  + \ldots + \dfrac{1}{a_n}\mu_n^N\right).
	\end{align} 
	Finally, observe that for all  $x\in N\setminus H$ and  $h\in H$, since $H\unlhd N$, $x^N$ and $(xh)^N$ intersect $H$ trivially. Therefore, $\mu_j^N(x) = 0 = \mu_j^N(xh)$  for all $j=1,\ldots,n$. Substituting into Equation \eqref{1}, we conclude that $\chi(x) = 0=\chi(xh)$, which completes the proof. 
\end{proof}

\end{document}
